\Subhead{4.4.1}{Оптические компрессоры и получение фемтосекундных импульсов}

Импульс описывается спектром $\widetilde A(\omega)=|\widetilde A|\exp[\ii\phi(\omega)]$. Постоянная фаза меняет фазу несущей, линейная — задержку, квадратичная — зависимость групповой задержки от частоты. Для гауссовского преобразовательно-ограниченного импульса по полным ширинам интенсивности $\Delta t\,\Delta\nu=2\ln2/\pi\simeq0.441$. Это минимальная длительность для \emph{данной гауссовой формы}, а не универсальная константа всех импульсов.

Компрессор вносит дисперсию групповой задержки, противоположную исходной: $\phi''_{\rm in}+\phi''_{\rm comp}\simeq0$. В паре дифракционных решёток частоты проходят разные оптические пути; призменные пары и чирпированные зеркала также управляют спектральной фазой. Важно компенсировать не только второй, но при широкой полосе и высшие порядки. Компрессор без расширения спектра не создаёт произвольно короткий импульс.

Самофазовая модуляция $\Delta\phi(t)=k_0n_2I(t)L$ расширяет спектр; последующая фазовая компенсация может сжать импульс. Для усиления больших энергий применяют растяжение, усиление и обратное сжатие (CPA): меньшая пиковая интенсивность в усилителе уменьшает нежелательные нелинейности и разрушение. Остаточная пространственно-временная связь и угловая дисперсия ухудшают фокусировку.

Длительность нельзя достоверно восстановить только из спектра интенсивности без фазы. Автокорреляция даёт косвенную оценку при предположении формы; FROG/SPIDER восстанавливают больше фазовой информации. Для нескольких оптических периодов существенна фаза между несущей и огибающей.
